% Artifact appendix for SCOPE, based on ae-template.tex (V20240722).
\documentclass{sigplanconf}

\usepackage{enumitem}
\usepackage{hyperref}

\begin{document}

\special{papersize=8.5in,11in}

\appendix
\section{Artifact Appendix}
\small

\subsection{Abstract}

The SCOPE artifact contains the software models, accuracy experiments, RTL,
synthesis reports, plotting scripts, and compact measurements needed to
evaluate the paper.  It reproduces Tables~3--5 and Figures~13--21, including
CPU-modeled accelerator performance, GPU model quality and nonlinear
approximation accuracy, and area/power results extracted from native synthesis
reports.  A hardware-free reviewer path regenerates the submitted evidence and
independently validates it.  Fresh CPU, GPU, RTL, and synthesis workflows are
also included.

\subsection{Artifact check-list (meta-information)}

\begin{itemize}[leftmargin=*]
  \item \textbf{Algorithms/programs:} The artifact includes SCOPE/SCNA
    approximation, LLMCompass, SCALE-Sim, OSTQuant, PyTorch, Chisel RTL, and
    evidence validators.
  \item \textbf{Compilation/transformations:} The Python and Cython programs
    run directly after installation.  The optional RTL workflow uses sbt,
    Scala, and Chisel to elaborate SystemVerilog.
  \item \textbf{Models/data:} The model set comprises OPT-6.7B, Llama-2-7B,
    Llama-3-8B, Qwen2.5-7B, and Qwen3-8B.  The data sets are WikiText-2, Pile
    validation, ARC-Easy, HellaSwag, PIQA, WinoGrande, and synthetic function
    samples.
  \item \textbf{Run-time environment:} The validated environment is 64-bit
    GNU/Linux with Python~3.12.  Fresh GPU execution additionally uses a CUDA
    12.8-compatible runtime and may use Slurm.
  \item \textbf{Hardware:} Bundled evidence and modeled performance run on a
    CPU.  Fresh quality experiments use NVIDIA GPUs.  Repeating synthesis
    requires the synthesis environment described below.
  \item \textbf{Execution/metrics:} Automated Make workflows report latency,
    throughput, speedup, perplexity, zero-shot accuracy, MSE, area, and power.
  \item \textbf{Output:} Outputs include CSV, JSON, and Markdown validation
    reports, PDF and PNG plots, tables, and generated SystemVerilog.
  \item \textbf{Experiments:} The artifact covers Tables~3--5,
    Figures~13--21, RTL regeneration, and the synthesis-report audit.
  \item \textbf{Resources:} Bundled review requires 2 CPU cores, 8~GB RAM,
    5~GB plus dependency caches, and about 15 seconds.  The fresh CPU suite
    requires 8 cores, 16~GB RAM, 20~GB, and about 37 minutes on dual
    EPYC~9654.  The fresh GPU suite takes about 5 hours with up to 16 H100 GPUs,
    excluding setup, downloads, and scheduler queue time.
  \item \textbf{Preparation:} Preparing the workflow takes about 10--30
    minutes with network access, excluding model approval and downloads.
  \item \textbf{Availability:} The source is available at
    \url{https://github.com/xsun2001/SCOPE-MICRO26-AE}.  The archived artifact
    is available at \url{https://doi.org/10.5281/zenodo.21472477}.
  \item \textbf{Licenses:} Component licenses are included.  LLMCompass uses
    BSD-3-Clause, OSTQuant uses Apache-2.0, and SCALE-Sim uses MIT.  External
    model and data licenses apply to downloads.
  \item \textbf{Automation:} The workflow uses GNU Make, Python, and optional
    Slurm job arrays.
\end{itemize}

\subsection{Description}

\subsubsection{How to access}

Clone or unpack the artifact above and enter its root directory.  The top-level
\texttt{README.md} is the reviewer entry point.  The file
\texttt{AE\_SUBMISSION.md} lists claims and resources, and every
\texttt{experiments/fig-*} or
\texttt{experiments/tbl-*} directory contains its own commands and expected
outputs.  Compact reference inputs and reproduction targets are in
\texttt{data} and \texttt{expected-results}.  Each \texttt{actual-results}
directory is ignored and reserved for fresh user executions.

\subsubsection{Hardware dependencies}

Bundled evidence needs no accelerator or proprietary software.  Fresh CPU
experiments need roughly 8 cores and 16~GB RAM.  Fresh Figure~16 and Table~5
need a CUDA GPU with at least 80~GB device memory.  Figures~17 and~20 need at
least 16~GB.  Table~5 additionally needs about 192~GB host memory and 200~GB
free storage.  Independent experiments may share up to 16 GPUs.

\subsubsection{Software dependencies}

Required tools are GNU Make, Bash, Python~3.10 or newer (3.12 validated), and
\texttt{uv} or \texttt{venv}/\texttt{pip}.  Fresh GPU runs use the pinned
\texttt{requirements/accuracy.txt} environment (PyTorch~2.10.0+cu128,
Transformers~4.51.0, and LM Evaluation Harness~0.4.4).  RTL regeneration uses
OpenJDK~21 and sbt.  Slurm is optional.

\subsubsection{Data sets and models}

CPU inputs and compact GPU evidence are included.  Fresh GPU runs download the
datasets listed above and require Hugging Face-format checkpoints under
\texttt{MODEL\_ROOT}.  Model weights and large Table~5 intermediates are not
redistributed.  Reviewers must obtain weights under the applicable terms.

\subsection{Installation}

From the artifact root, run:
\begin{verbatim}
make setup
make evidence
make validate
\end{verbatim}
The first command creates \texttt{.venv} and installs dependencies.  An
existing environment can instead be selected with
\texttt{PYTHON=/path/to/python}.
Network access is normally needed for initial package, model, and data fetches.

\subsection{Experiment workflow}

\begin{itemize}[leftmargin=*]
  \item \textbf{Fast review:} \texttt{make evidence \&\& make validate}
    regenerates compact tables/plots in ignored run directories and validates
    all packaged CPU/GPU results without changing tracked evidence.
  \item \textbf{Fresh CPU/RTL:} \texttt{make reproduce-cpu JOBS=8} runs the
    modeled performance suite, hardware extraction, and Chisel elaboration.
  \item \textbf{Fresh GPU:} copy \texttt{config/local.env.example} to
    \texttt{config/local.env}, set \texttt{MODEL\_ROOT}, then run
    \texttt{make reproduce-gpu}.  Use \texttt{EXECUTOR=local WORKERS=1} on an
    allocated GPU or \texttt{EXECUTOR=slurm WORKERS=N} for job arrays.
\end{itemize}

\subsection{Evaluation and expected results}

A successful packaged review prints \texttt{status: pass} for the following
checks:
\begin{itemize}[leftmargin=*]
  \item 68/68 CPU checks and 36 Table~3 groups recomputed from 720 raw H100
    samples.
  \item 80/80 Figure~16, 20/20 Table~5, 36/36 Figure~17, and 18/18 Figure~20
    comparisons.
  \item All 11 SCNA rows in Table~4 and zero absolute author paths.
\end{itemize}
Figures~13--15 and~21 match the reported performance and speedup targets.
Figures~18--19 are regenerated from 112 filtered native synthesis report sets.
Figure~16 and Table~5 both use WikiText-2 perplexity and the unweighted mean
accuracy over ARC-Easy, HellaSwag, PIQA, and WinoGrande.  Table~4 reproduces
SCOPE/SCNA accuracy.  Its other columns are literature reference values.

\subsection{Experiment customization}

Use \texttt{JOBS=N} for CPU workers and \texttt{WORKERS=N} for GPU job-array
width.  Make's \texttt{-j} runs independent experiment targets concurrently.
The product of concurrent targets and \texttt{WORKERS} must remain within the
GPU allocation.
Individual targets (for example, \texttt{make fig-16} or \texttt{make tbl-5})
permit partial reproduction.  Table~5 checkpoints may be reused through
\texttt{TABLE5\_CHECKPOINT\_SOURCE}.

\subsection{Notes}

The B200/B300/H100/AWSv4/TPUv6e performance claims are software-model results
and do not require those devices.  Archived synthesis used Design Compiler
V-2023.12 and TSMC 28~nm libraries at 1~GHz.  These proprietary inputs are
needed only to repeat synthesis, not to audit the bundled native reports.
Provenance placeholders are documented in \texttt{PROVENANCE.md}.

\subsection{Methodology}

Artifact reviewing and badging follow
\url{https://www.acm.org/publications/policies/artifact-review-and-badging-current}
and \url{https://cTuning.org/ae}.

\end{document}
